\chapter{Modeling for Operational Databases}

This chapter discusses modeling methods and techniques for operational databases, introducing the main notation elements of the most common models.

\section{Information Modeling}

When a database system is created, it is based on some model that reproduces the essential characteristics of the real world information it is intended to represent. In general, models can be either \textit{iconic}, if they retain physical characteristics of the entities they represent, or \textit{symbolic}, if they use symbols to represent entities and their relationships. In the context of information systems, symbolic models are used.
\BoxDef{Symbolic Model}{
    A \textbf{symbolic model} is a subjective, formal representation of ideas and knowledge about some aspects of the real world, called domain of discourse, designed to serve a specific purpose.
}
A model simplifies reality by focusing on what is important for the task at hand, and uses a formal language to express this simplified representation.\\
Modeling in databases is broken down into three levels of abstraction:
\begin{itemize}
    \item \textbf{Conceptual level}, in which the basic structure of the database is defined at an high-level, after analyzing the problem and considering any user requirement. Some examples of conceptual models are the \textbf{Entity-Relationship} (\textbf{ER}) model and the \textbf{Object Data Model} (\textbf{ODM});
    \item \textbf{Logical level}, where the conceptual model is mapped to a more detailed description of the data structures, including all attributes and constraints. The model is still independet from any specific DBMS. An example of logical model is the \textbf{Relational Data Model};
    \item \textbf{Physical level}, where the physical data structures are defined for the elements in the logical model. The choices done at this level must take in consideration the features of the DBMS used to implement the database.
\end{itemize}

\subsection{What to Model}
The reality to be modeled is considered in terms of concrete knowledge, abstract knowledge, procedural knowledge, dynamics, and communications.

\paragraph{Concrete Knowledge} Concrete knowledge refers to specific facts known of the reality to be represented. We can assume reality consists of \textit{entities}, \textit{properties}, and \textit{relationships}.
\BoxDef{Entity, Property, Relationship}{
    An \textbf{entity} is anything for which certain facts should be recorded, independently of the existence of other entities.

    A \textbf{property} is a fact about an entity which is not meaningful in itself, but only because it describes an entity of interest.

    A \textbf{relationship} is a fact which correlates independent entities.
}
For example, we want to model the administration system of a library. Examples of entities can be a book, a bibliographic description, a registered user, or a loan record. Properties may be a user's name and address, or a book's title and author. Entities with the same properties are said to have the same \textbf{type}, and are classified in the same \textbf{collection}, or \textbf{entity set}. The books titled ``1984'', ``Moby Dick'', and ``The Great Gatsby'' are all part of the same entity set, and are all of type \textit{book}. As for relationships, examples could be the relationship between a bibliographic description and one or more books (the same book may be owned multiple times by the library), or between a user and their loan records.

\paragraph{Abstract Knowledge} Abstract knowledge concerns facts which impose certain restrictions on admissible values of concrete knowledge and on the way these values can evolve over time, or expresses rules on how to derive new information. These rules are called \textbf{integrity constraints}. In the library example, abstract knowledge could include rules about the data type of attributes, such that book properties \textit{title} and \textit{author} must be strings, while \textit{length} must be a non-negative integer; another example may be the rule that a loan can only be borrowed for two weeks at a time, and must be associated with the \textit{id} of a registered user; the age of a user can be automatically calculated by subtracting the \textit{birthdate} from the current date. Relationships can also have constraints that limit the possible correlated entities. The most important ones are \textbf{structural constraints}, which are:
\begin{itemize}
    \item \textbf{Cardinality}, \textit{one} or \textit{many}, to specify how many entities of one collection can be associated with entities of another collection, e.g., a bibliographic description can be associated with \textit{many} book copies, while a book copy is associated with only \textit{one} bibliographic description, so the relationship is called \textit{one-to-many};
    \item \textbf{Participation}, \textit{total} or \textit{partial}, to specify whether an entity of one collection can have entities of another collection associated with it, e.g., a loan record must be associated with one registered user, so the participation from the loan record side is \textit{total}; on the other hand, a registered user may have never loaned a book, so the participation from the registered user side is \textit{partial}. 
\end{itemize}

\section{Data Models}

To construct the conceptual model for an operational database, we define a \textbf{schema}, a collection of time-invariant definitions which model the structure of admissible data (including integrity constraints), and procedural knowledge (i.e., the elementary operations applied to concrete knowledge to cause changes). Different formalisms, each with a specific data model, can be used to define the schema.
\BoxDef{Data Model}{
    A \textbf{Data Model} is a set of abstraction mechanisms, with associated operators and implicit integrity constraints, used to define a database schema.
}
For operational databases, we will see the Object Data Model and the Relational Data Model, for the conceptual and logical levels respectively.

\subsection{Object Data Model}
The basic mechanisms of the ODM are:
\begin{itemize}
    \item \textbf{Objects}: they are the equivalent of entities. An object is a software entity with an internal state represented by instance variables, and a behaviour described by a set of methods. Each object has its own identity that persists over time, even if its internal state change;
    \item \textbf{Object types}: just like an entity belongs to a type, an object belongs to an object type. An object type describes the state fields and the implementation of methods common to all objects of that type;
    \item \textbf{Classes}: a modifiable sequence of objects with the same type. They are the equivalent of entity sets. Each class introduces the definition of a type and a constructor for values of this type, as well as a name to denote the class itself;
    \item \textbf{Relationships}: classes have relationships with other classes, which are characterized by a cardinality and partiality. A relationship may also have attributes, or be recursive (i.e., a class has a relationship with itself);
    \item \textbf{Inheritance}: a mechanism that allows something to be defined by only describing how it differs from a previously defined one. It is distinct from subtyping, which is a relation between types such that when one type $A$ is a subtype of another type $B$, then any operation applicable to type $B$ can also applied to type $A$ (although inheritance is one way to define subtypes). A type can be defined from a single supertype (\textit{single inheritance}) or from several supertypes (\textit{multiple inheritance}).
    \item \textbf{Class hierarchies}: similarly to types, classes can also be organized in hierarchies. If $C_1$ is a subset of $C_2$, then $C_1$ is a \textit{subclass} of $C_2$, and this implies that:
        \begin{itemize}[noitemsep]
            \item The type of elements in $C_1$ is a subtype of the type of the elements in $C_2$;
            \item The elements in $C_1$ are a subset of the elements in $C_2$, and $C_1$ inherits any relationship defined on $C_2$.
        \end{itemize}
    When the same superclass is shared by two or more subclasses, two constraint can be defined: \textbf{overlap} and \textbf{covering}. A no overlap constraint specifies that two subclasses cannot share any element. In the absence of this constraint, the same object may belong to more that one subclass. A covering constraint specifies that the objects in the subclasses must collectively include all elements in the superclass. In the absence of this constraing, the superclass may have elements that do not belong to any subclass. When both constraints are present, we call the hierarchy a \textbf{generalization}.
\end{itemize}
The following figures illustrate a possible graphical notation used for ODM. There is no standard notation; the one used here is based on the notation used in UML.
\begin{figure}[H]
    \centering
    \includesvg[width=0.55\textwidth]{img/ODM_objects.svg}
    \caption{Graphical notation for objects, object types, and classes. Classes can be shown with different levels of detail.}
\end{figure}
\begin{figure}[H]
    \centering
    \includesvg[width=0.45\textwidth]{img/ODM_relationships_1.svg}
    \caption{Simple relationships between classes. Each relationship is represented by an arc with a label. Multivalued relationships are represented with a double arrow. Optional relationships have a crossed line. Recursive relationships require labels on the ends of the arcs to clarify the meaning.}
\end{figure}
\begin{figure}[H]
    \centering
    \includesvg[width=0.5\textwidth]{img/ODM_relationships_2.svg}
    \caption{Relationships with attributes.}
\end{figure}
\begin{figure}[H]
    \centering
    \includesvg[width=0.6\textwidth]{img/ODM_hierarchies.svg}
    \caption{Class hierarchies.}
\end{figure}

\subsection{Relational Data Model}

The Relational Data Model describes databases in terms of sets of \textbf{tuples} (also called \textbf{records}), and associations among elements are expressed in terms of values of their attributes.
\BoxDef{Relational Database}{
    A \textbf{relational database} is described by a set of relation schemas $R: \{ T \}$ defined as follows:
    \begin{itemize}[noitemsep]
        \item \textbf{integers}, \textbf{floats}, \textbf{booleans}, and \textbf{strings} are primitive types.
        \item If $T_1, \ldots, T_n$ are primitive types, and $A_1, \ldots, A_n$ are distinct attribute names, then $T = (A_1:T_1, \ldots, A_n:T_n)$ is a \textbf{tuple type} of degree $n$. The attribute order is unimportant.
        \item If $T$ is a tuple type, then $\{T\}$ is a \textbf{relation type}.
        \item A \textbf{relation schema} $R : \{T\}$ is a variable $R$ with a relation type.
    \end{itemize}
}
For brevity, instead of writing $R:\{T\}$, we will use $R(T)$. Also, when the type of the attributes in a schema is unimportant, we will simply write it as $R(A_1, \ldots, A_n)$.

\BoxDef{Tuple}{
    A \textbf{tuple} $(A_1 := V_1, \ldots, A_n := V_n)$ of type $T = (A_1:T_1, \ldots, A_n:T_n)$ is a set of pairs $(A_i, V_i)$ with $V_i$ of primitive type $T_i$.
}

\BoxDef{Relation}{
    A \textbf{relation} (or schema instance) is a finite set of tuples with the same type $T$.
}
A relation does not contain any duplicate tuples, since it is a set and not a multi-set. The order of relation elements and attributes is unimportant.\\
The following statements hold true:
\begin{itemize}[noitemsep]
    \item Two tuple types are equal if and only if they have the same degree, same attributes, and same type of the attributes with the same name.
    \item Two relation types are equal if and only if they have the same tuple types.
    \item Two tuples of the same type are equal if and only if they have the same set of attribute-value pairs. 
    \item Two relations of the same type are equal if and only if they have equal tuples.
\end{itemize}

\BoxDef{Key}{
    A \textbf{key} for a relation is the minimal subset of attributes whose values identify a tuple. Out of all the possible keys, the database designer chooses a \textbf{primary key} (usually the shortest).
}
Relationships between relations are expressed using \textbf{foreign keys}, which take as values those of the primary key of another relation. For example, consider the two relations:
\begin{center}
    \textit{\textbf{Students}(\underline{StudentCode}, Name, City, BirthYear)}\\
    \textit{\textbf{Exams}(\underline{Subject}, \underline{Candidate}, Date, Grade)}
\end{center}
where attribute \textit{StudentCode} is the primary key of \textit{Students}, and the pair \textit{Subject}, \textit{Candidate} is the primary key of \textit{Exams}. Since the attribute \textit{Candidate} takes values that match those of \textit{StudentCode}, it is a foreign key.

\paragraph{Relational Algebra} Relational algebra is a set of operators on relations whose results are also relations. Expressions in relational algebra are called \textbf{queries}, and the standard language used by DBMSs to write queries is SQL (Structured Query Language): it is a declarative language, meaning that the user specifies what data to retrieve, and not how to retrieve it. There are six fundamental operations, plus additional operations that can be derived from them. These operations are:
\begin{itemize}
    \item \textbf{Rename} $\rho_{A_1 \leftarrow B_1, \ldots, A_m \leftarrow B_m} (R)$: assigns a new name to attributes. $A_1, \ldots, A_m$ are attributes of relation $R$, and $B_1, \ldots, B_m$ are not. The result is a relation with type $\{(B_1:T_1, \ldots, B_m:T_m)\}$, whose tuples are those of $R$ with renamed attributes.
    \item \textbf{Project} $\pi_{A_1, \ldots, A_m} (R)$: selects a subset of attributes from the relation $R$, deleting any duplicate. The result is a relation with type $\{(A_1:T_1, \ldots, A_m:T_m)\}$, whose tuples are those of $R$ restricted to attributes $A_1, \ldots, A_m$.
    \item \textbf{Generalized project} $\pi_{\textit{exp}_1 \text{ AS }B_1, \ldots, \textit{exp}_n \text{ AS }B_n} (R)$: a combination of projection and renaming that also allows the output to contain the result of expressions. $\textit{exp}_1,\ldots,\textit{exp}_n$ are arithmetic expressions that involve constants and attributes of $R$, and $B_1, \ldots, B_n$ are attribute names. The result is a relation with type $\{(B_1:T_{\text{exp}_1}, \ldots, B_n:T_{\text{exp}_n})\}$.
    \item \textbf{Select} $\sigma_{\textit{cond}} (R)$: selects the tuples in $R$ that satisfy the condition $\textit{cond}$. This condition is a propositional logic expression ($\land, \lor, \lnot$) over comparison predicates ($\leq, <, \geq, >, =, \not =$) and arithmetic expressions ($+,-,*,\div$).
    \item \textbf{Grouping} $_{A_1, \ldots, A_n} \gamma_{f_1, \ldots, f_k} (R)$: returns a relation whose tuples are grouped by the values of the specified attributes, and applies the specified aggregate functions to the remaining attributes. $A_1, \ldots, A_n$ must be attributes of $R$, and $f_1, \ldots, f_k$ are aggregate functions (min, max, count, sum, average, etc.) calculated over one or more of the grouped attributes. The relation type of the result is $\{(A_1:T_1, \ldots, A_n:T_n, f_1:T_{f_1}, \ldots, f_k:T_{f_k})\}$.
    \item \textbf{Generalized grouping} $_{A_1, \ldots, A_n} \gamma_{f_1 \text{ AS } F_1, \ldots, f_k \text{ AS } F_k} (R)$: extends the grouping operation by allowing the renaming of attributes obtained via aggregation functions. The relation type of the result is $\{(A_1:T_1, \ldots, A_n:T_n, F_1:T_{f_1}, \ldots, F_k:T_{f_k})\}$.
    \item \textbf{Union} $R \cup S$: returns the union of the two sets of tuples in $R$ and $S$, given that they have the same relation type.
    \item \textbf{Set difference} $R - S$: returns the set of tuples in $R$ that are not also in $S$, given that they have the same relation type.
    \item \textbf{Product} $R \times S$: given $R$ of type $\{(A_1:T_1, \ldots, A_n:T_n)\}$ and $S$ of type $\{(A_{n+1}:T_{n+1}, \ldots, A_{n+m}:T_{n+m})\}$ with disjoint set of attributes, the product operator returns a relation of type $\{A_1:T_1, \ldots, A_n:T_n, A_{n+1}:T_{n+1}, \ldots A_{n+m}:T_{n+m}\}$, whose tuples are all possible combinations of tuples in $R$ and $S$.
    \item \textbf{Join} $R \underset{A_i = A_j}{\bowtie} S$: combines tuples from $R$ and $S$ that have the same values for the attributes $A_i$ of $R$ and $A_j$ of $S$. The two relations have type $\{(A_1:T_1, \ldots, A_n:T_n)\}$ and $\{(A_{n+1}:T_{n+1}, \ldots, A_{n+m}:T_{n+m})\}$, and have a disjoint set of attributes. This operation is equivalent to calculating $\sigma_{A_i = A_j} (R \times S)$.
    \item \textbf{Intersection} $R \cap S$: returns the set of tuples that are present in both $R$ and $S$, given that they have the same relation type.
    \item \textbf{Natural join} $R \bowtie S$: a special case of join that is only applicable when the two relations have attributes with the same name. It combines tuples in $R$ and $S$ that have the same values for all attributes with the same name in both.
    \item \textbf{Division} $R \div S$: let $XY$ be the attributes of $R$, and $Y$ the attributes of $S$. Then, the result of the division between $R$ and $S$ is a relation with attributes $X$ such that
        \begin{equation*}
            W = \{ w | \forall t \in S.(w \circ t \in R) \}
        \end{equation*}
    where $\circ$ is the concatenation operator. In other words, it returns all the values of attributes $X$ in $R$ which appear in a tuple with all values of attributes $Y$ in $S$. \\
    If $S \times W = R$, then $R \div S = W$.
\end{itemize}
There are also a series of \textbf{equivalence rules} to simplify expressions:
\begin{align*}
    &\pi_A(\pi_{A,B}(R)) \equiv \pi_A(R) &(\text{Cascading of projections})\\
    &\sigma_{\textit{cond}_1}(\sigma_{\textit{cond}_2}(R)) \equiv \sigma_{\textit{cond}_1 \land \textit{cond}_2}(R) &(\text{Cascading of selections})\\
    &\sigma_{\textit{cond}_R \land \textit{cond}_S}(R \bowtie S) \equiv \sigma_{\textit{cond}_R}(R) \bowtie \sigma_{\textit{cond}_S}(S) &(\text{Commutativity of selection and join})\\
    &R \bowtie (S \bowtie T) \equiv (R \bowtie S) \bowtie T &(\text{Associativity of join})\\
    &R \bowtie S \equiv S \bowtie R &(\text{Commutativity of join})\\
    &\sigma_{\textit{cond}_X}(_X \gamma_{F}(R)) \equiv _X \gamma_{F}(\sigma_{\textit{cond}_X}(R)) &(\text{Selection distributes over grouping})
\end{align*}